\subsubsection{Общая архитектура}

Для проведения эксперимента была реализована программная система из нескольких
контейнеров. Пользователь работает с Vue-приложением, выбирает модель хранения,
профиль данных и сценарий нагрузки. Клиентское приложение отправляет запросы в
Express API, а сервер перенаправляет их в нужный экземпляр сервиса генерации нагрузки (loadgen).

В системе постоянно запущены два loadgen-сервиса. Первый подключается к
PostgreSQL и выполняет прогоны для моделей \texttt{pg-jsonb} и
\texttt{pg-normalized}. Второй подключается к MongoDB и выполняет прогоны для
\texttt{mongo-nested} и \texttt{mongo-normalized}. Такой вариант позволяет
использовать общий код генератора нагрузки, но разделить подключение к разным
СУБД.

Метрики собираются через Prometheus. Loadgen публикует собственные метрики
операций, PostgreSQL exporter и MongoDB exporter показывают состояние баз
данных, а Grafana отображает итоговые графики. Архитектура приложения показана
на рисунке.

\begin{figure}
    \centering
    \includegraphics[
        width=\textwidth,
        height=1\textheight,
        keepaspectratio
    ]{architecture.pdf}
    \caption{Архитектура веб-приложения}
\end{figure}

\FloatBarrier
